
\section{1. Introduction}

\subsection{1.1 Motivation}

The softmax attention mechanism (Vaswani et al., 2017) scales quadratically 
with sequence length $L$:

\[
\text{Attn}(Q, K, V) = \text{softmax}\!\left(\frac{Q K^\top}{\sqrt{d_h}}\right) V
\tag{1.1}
\]

producing an $L \times L$ attention matrix that requires $\mathcal{O}(L^2)$ 
memory and time. For a sequence of 4096 tokens at 12 heads, the attention 
matrix alone occupies \textbf{200 MB at FP16}, dwarfing the model's own weight 
footprint.

\textbf{Low-rank approximation} reduces this to $\mathcal{O}(L m)$ by projecting 
onto $m \ll L$ landmarks. The Nystr m method (Williams \& Seeger, 2001) does 
this by:

\[
K \approx \tilde{K} = C W^{-1} C^\top
\tag{1.2}
\]

where $C \in \mathbb{R}^{L \times m}$ is the query-landmark kernel and 
$W \in \mathbb{R}^{m \times m}$ is the landmark-landmark kernel. The 
Nystr mformer (Xiong et al., 2021) applies this to attention.

\textbf{The problem with standard Nystr m.} It assumes a symmetric kernel, 
requiring the same landmarks for queries and keys. In attention, however, 
queries and keys play different roles: queries ask "what am I looking for?"; 
keys answer "what am I?" Using the same landmarks ignores this asymmetry.

\subsection{1.2 Problem Statement}

Given query $Q \in \mathbb{R}^{L \times d}$, key $K \in \mathbb{R}^{L \times d}$, 
and value $V \in \mathbb{R}^{L \times d}$, we seek a low-rank approximation 
of the attention output:

\[
\text{Attn}(Q, K, V) \approx \text{Attn}_\text{low-rank}(Q, K, V)
\tag{1.3}
\]

such that:

1. \textbf{Complexity:} $\mathcal{O}(L m)$ instead of $\mathcal{O}(L^2)$, with $m \ll L$.
2. \textbf{Asymmetry:} Independent landmark selection for $Q$ and $K$.
3. \textbf{Differentiability:} End-to-end training.
4. \textbf{Compatibility:} Works with pretrained models without retraining.
5. \textbf{Stability:} Numerically robust for FP16 and long sequences.

\subsection{1.3 Our Approach: Primal Attention}

We propose \textbf{Primal Attention}, which decomposes the attention matrix 
asymmetrically:

\[
A \approx F_1 F_2^{-1} F_3
\tag{1.4}
\]

where:

\begin{itemize}
\item $F_1 = Q K_\text{land}^\top \in \mathbb{R}^{L \times m}$: query vs landmark keys
\item $F_2 = Q_\text{land} K_\text{land}^\top \in \mathbb{R}^{m \times m}$: landmark-landmark
\item $F_3 = Q_\text{land} K^\top \in \mathbb{R}^{m \times L}$: landmark queries vs keys
\end{itemize}

This is the \textbf{primal representation} of the asymmetric kernel (in contrast 
to the dual representation which uses $L \times L$ kernels). The pseudoinverse 
$F_2^{-1}$ is computed via \textbf{Schulz iteration}:

\[
X_{t+1} = X_t (2I - F_2 X_t)
\tag{1.5}
\]

which converges quadratically to $F_2^{-1}$ in a few steps and is differentiable.

\textbf{Landmark selection.} We use \textbf{segment pooling}: divide the sequence into 
$m$ equal chunks and average each chunk. This gives uniform coverage 
regardless of content.

\subsection{1.4 Contributions}

\begin{itemize}
\item \textbf{An asymmetric kernel SVD formulation.} Primal Attention allows independent
\end{itemize}
  landmarks for queries and keys, generalizing standard Nystr m approximation.
\begin{itemize}
\item \textbf{A differentiable pseudoinverse.} Schulz iteration provides end-to-end
\end{itemize}
  trainable pseudoinverse without explicit eigendecomposition.
\begin{itemize}
\item \textbf{A production-hardened implementation.} Supports flat, nested, and causal
\end{itemize}
  architectures with GQA, RoPE, and safe injection.
\begin{itemize}
\item \textbf{A systematic evaluation.} Benchmarked on DistilBERT, RoBERTa, and GPT-2
\end{itemize}
  on SST-2 with ablations over landmark count and pseudoinverse iterations.
\begin{itemize}
\item \textbf{Full reproducibility.} Code, checkpoints, and evaluation scripts released
\end{itemize}
  under MIT.

\subsection{1.5 Paper Organization}

Section 2 reviews related work. Section 3 provides background on softmax 
attention, Nystr m approximation, and pseudoinverse. Section 4 presents 
Primal Attention. Section 5 describes integration with four architectures. 
Section 6 details the experimental setup. Section 7 reports results. Section 
8 analyzes where the method succeeds and fails. Section 9 discusses practical 
guidance. Section 10 lists limitations. Section 11 concludes.



\subsection{Notation}
